\documentclass[10pt,a4paper]{amsart}
\usepackage[margin=19mm]{geometry}
\usepackage{amsmath,amssymb,mathtools}
\usepackage[scr=rsfs,cal=euler]{mathalfa}
\usepackage{xfrac}
\usepackage[unicode,hidelinks,bookmarks=false]{hyperref}
\newcommand{\ie}{i.e.\ }
\newcommand{\cf}{cf.\ }
\newcommand{\resp}{respectively\ }
\newcommand{\Subsection}{\subsection}
\providecommand{\nobreakdash}{\nobreak\hbox{-}}
\setlength{\emergencystretch}{2em}
\multlinegap0pt
\usepackage{bm}
\usepackage{bbm}
\usepackage{dsfont}
\usepackage{xspace}
\usepackage{cancel}
\usepackage{mathtools}
\usepackage{amsthm}
\makeatletter
\@ifundefined{definition}{\newtheorem{definition}{Definition}}{}
\@ifundefined{lemma}{\newtheorem{lemma}{Lemma}}{}
\@ifundefined{observation}{\newtheorem{observation}{Observation}}{}
\makeatother
\newcommand*{\eqcommentnear}[1]{&&\text{[#1]}}
\newcommand*{\eqcomment}[2][\qquad]{#1 &&\text{[#2]}}
\newcommand{\Ssm}{S_{sm}}
\newcommand{\divides}{\mid}
\newcommand{\ltdivides}{\mid_{lt}}
\newcommand{\ltfrac}[2]{\frac{\,#1\,}{\,#2\,}{}_{{\scriptscriptstyle lt}}}
\newcommand{\ltsm}{\text{LT}_{sm}}
\newcommand{\piltd}{\pi_{lt}}
\newcommand{\tail}{\text{tail}_{sm}}
\title[Gr\"obner Bases Native to Term-ordered Commutative Algebras,]{Gr\"obner Bases Native to Term-ordered Commutative Algebras, with Application to the Hodge Algebra of Minors}
\author{Joshua A. Grochow, Abhiram Natarajan}
\date{Source: arXiv:2510.11212v1 (CC BY 4.0)}
\begin{document}
\maketitle

\noindent\emph{Source and licence.} Gr\"obner Bases Native to Term-ordered Commutative Algebras, with Application to the Hodge Algebra of Minors. Version arXiv:2510.11212v1 (CC BY 4.0), distributed under the Creative Commons Attribution 4.0 International licence, as recorded in the arXiv licence metadata for this exact version and shown on the abstract page https://arxiv.org/abs/2510.11212v1. Full rights notice: licenses/arxiv-2510.11212-CC-BY-4.0.txt.\par\smallskip

\noindent\textit{Editorial scope.} Counted unit: the Lemma whose source label is \texttt{lem:Ssm-prod} in the article (source0.tex lines 1573--1579, source label \texttt{lem:Ssm-prod}), delivered together with the article's own complete original proof of it (source0.tex lines 1581--1608, one \texttt{proof} environment of 28 lines). No mathematics is written, repaired or re-derived: statement and proof are the article's own text line for line. Required context retained uncounted: lem:mon-div (lines 705--709); obs:mlt (lines 1246--1252); def:ltdivides-modules (lines 1322--1324). Every definition, display and label of those lines is retained unchanged. Omitted from this package and not redistributed: 1--704, 710--1245, 1253--1321, 1325--1572, 1580--1580, 1609--3283. Nothing inside a retained excerpt is omitted or altered: every retained body is delivered exactly as the article's lines are written, so no omission receipt of any kind accompanies this package. Citations: the retained excerpts carry no citation command at all, so no bibliography entry of the article is transcribed and no reference is redistributed. Reference adaptation: none was needed and none was made; every reference inside the delivered unit resolves inside it, so no unresolved reference or citation is produced. Printed slips kept literal: no printed word, symbol, hypothesis or number of the counted statement or of its proof was altered. Layout and class: the article's own class is \texttt{amsart} with packages including amsmath, amssymb, amsthm, xspace, hyperref, mathtools, geometry, xcolor, tabmac, float, enumitem, nicefrac, parskip, algorithm; the wrapper is the lane's pinned \texttt{amsart} editorial wrapper at 10pt on A4, which restates the theorem-like environments the retained text needs and adds the article's own macro definitions that the retained text needs (\texttt{\textbackslash Ssm}, \texttt{\textbackslash divides}, \texttt{\textbackslash eqcomment}, \texttt{\textbackslash eqcommentnear}, \texttt{\textbackslash ltdivides}, \texttt{\textbackslash ltfrac}, \texttt{\textbackslash ltsm}, \texttt{\textbackslash piltd}, \texttt{\textbackslash tail}), with extra wrapper packages (bm, bbm, dsfont, xspace, cancel, mathtools). The delivered numbering is the unit's own. Assets: no retained span contains a figure environment, a graphics-inclusion command, a PDF-page inclusion, a table or a TikZ picture; no image file is copied into this package, the package declares no asset dependency and no image file is redistributed with it. Licence: version arXiv:2510.11212v1 of this article is distributed under the Creative Commons Attribution 4.0 International licence, as recorded in the arXiv licence metadata for this exact version and shown on the abstract page https://arxiv.org/abs/2510.11212v1. A full rights notice is delivered at licenses/arxiv-2510.11212-CC-BY-4.0.txt. Characters: every retained excerpt is pure ASCII, so the delivered engine, XeLaTeX with its default Latin Modern text font, reports no missing character.
\begin{lemma} \label{lem:mon-div}
Let $A'$ be a monomial pseudo-ASL. If $t,t'$ are standard terms such that $t \mid t'$, then there is a standard term $t''$ such that $tt'' = t'$.

If, furthermore, $A'$ admits a pseudo-ASL term order $\preceq$, then the standard term as above is unique.
\end{lemma}

\begin{observation}[store=obsmlt, note={Module of leading terms}]
\label{obs:mlt}
Given $(A^d, \preceq, A_{lt}^d)$, by slight abuse of notation let $\piltd \colon A^d \to A_{lt}^d$ be the map that applies $\pi_{lt} \colon A \to A_{lt}$ coordinate-wise. Then $\piltd$ is an $R$-linear bijection that restricts to the identity map on standard monomials, and for all standard monomials $m \in A, m' e_i \in A^d$, we have
\[
\pi_{lt}(m) \piltd(m' e_i) \in \{0, \piltd(\ltsm(mm' e_i))\}.
\]
\end{observation}

\begin{definition}[Divisibility of leading terms in modules] \label{def:ltdivides-modules}
Given $(A^d, \preceq, A_{lt}^d)$ and $f,g \in A^d$. If $\piltd(f) \divides \piltd(g)$, we write $f \ltdivides g$. If $m,m' \in A^d$ are standard terms such that $m \ltdivides m'$, we may further write $\ltfrac{m'}{m}$ for the unique (by the module analogue of Lemma~\ref{lem:mon-div}) standard term $t \in A_{lt}$ such that $\piltd(m)t = \piltd(m')$. 
\end{definition}

\begin{lemma} \label{lem:Ssm-prod}
Let $g,g' \in A^d$, $l$ a standard monomial in $\langle\piltd(\ltsm(g))\rangle \cap \langle \piltd(\ltsm(g'))\rangle$, and $\mu \in A$ a standard monomial that is compatible with $l$. Then there exist $h,h' \in A$ such that
\[
\mu \Ssm^{l}(g,g') = \Ssm^{\pi_{lt}(\mu) l}(g,g') + hg + h'g'
\]
with $\ltsm(hg)$ and $\ltsm(h'g')$ both strictly less than $\ltsm(\mu \pi_{lt}^{-1}(l))$.
\end{lemma}

\begin{proof}
By definition, we have
\begin{align*}
\mu \Ssm^{l}(g, g') & = \mu\left(\pi_{lt}^{-1}\left(\ltfrac{l}{\ltsm(g)}\right) g - \pi_{lt}^{-1}\left(\ltfrac{l}{\ltsm(g')}\right) g' \right) \\
 \end{align*}
 At this point, since we will use them frequently throughout this derivation, we introduce the shorter notation $\alpha := \pi_{lt}^{-1}\left(\ltfrac{l}{\ltsm(g)}\right)$ and $\beta := \pi_{lt}^{-1}\left(\ltfrac{l}{\ltsm(g')}\right)$. Continuing with this new notation in hand, we have
\begin{align}
 \mu \Ssm^{l}(g, g') & = \mu \alpha g - \mu \beta g' \nonumber \\
  & = \left(\ltsm(\mu \alpha) + \tail(\mu \alpha) \right) g - \left(\ltsm(\mu \beta) + \tail(\mu \beta) \right) g' \nonumber \\
  & = \left(\ltsm(\mu \alpha) g - \ltsm(\mu \beta) g'\right) + \left(\tail(\mu \alpha) g - \tail(\mu \beta) g'\right)  \label{eq:mussm}\\
\end{align}
We claim that the first summand here is $\Ssm^{\pi_{lt}(\mu) l}(g,g')$, and that we may take $h = \tail(\mu \alpha)$ and $h' = -\tail(\mu \beta)$ to satisfy the conclusion of the lemma. To prove both of these, we now analyze $\ltsm(\mu \alpha)$ and $\ltsm(\mu \beta)$.

We have
\begin{align*}
\pi_{lt}\left(\ltsm(\mu \alpha)\right) & = \pi_{lt}\left(\ltsm\left(\mu \pi_{lt}^{-1}\left(\ltfrac{l}{\ltsm(g)} \right) \right)\right) \\
 & = \pi_{lt}\left( \mu \right) \pi_{lt}\left( \pi_{lt}^{-1}\left(\ltfrac{l}{\ltsm(g)} \right)\right)  \eqcommentnear{by Obs.~\ref{obs:mlt}} \\
 & = \pi_{lt}(\mu) \cdot \ltfrac{l}{\ltsm(g)} \\
 & = \pi_{lt}(\mu) \frac{l}{\pi_{lt}(\ltsm(g))} \eqcomment{by Def.~\ref{def:ltdivides-modules}} \\
& = \frac{\pi_{lt}(\mu) l}{\pi_{lt}(\ltsm(g))} 
\end{align*}

It follows that $\ltsm(\ltsm(\mu \alpha) g) = \pi_{lt}(\mu) l$. A similar calculation yields that $\ltsm(\ltsm(\mu \beta) g') = \pi_{lt}(\mu) l$. It follows that
\[
\left(\ltsm(\mu \alpha) g - \ltsm(\mu \beta) g'\right) = \Ssm^{\pi_{lt}(\mu) l}(g,g')
\]
and $\ltsm(\tail(\mu \alpha) g)$ and $\ltsm(\tail(\mu \beta) g')$ are both strictly less than $\ltsm(\mu \pi_{lt}^{-1}(l))$. This completes the proof of the lemma.
\end{proof}

\end{document}
